\documentclass[10pt,a4paper]{amsart}
\usepackage[margin=19mm]{geometry}
\usepackage{amsmath,amssymb,mathtools}
\usepackage[scr=rsfs,cal=euler]{mathalfa}
\usepackage{xfrac}
\usepackage[unicode,hidelinks,bookmarks=false]{hyperref}
\newcommand{\ie}{i.e.\ }
\newcommand{\cf}{cf.\ }
\newcommand{\resp}{respectively\ }
\newcommand{\Subsection}{\subsection}
\providecommand{\nobreakdash}{\nobreak\hbox{-}}
\setlength{\emergencystretch}{2em}
\multlinegap0pt
\usepackage{amssymb}
\usepackage{mathtools}
\newtheorem{theorem}{Theorem}
\newtheorem{lemma}{Lemma}
\newcommand{\R}{\mathbb{R}}
\newcommand{\gs}{\gamma_{*}}
\newcommand{\mub}{\bar\mu}
\title{Radial solutions of quasilinear Hardy--Emden--Fowler equations: exact decay rates, Delaunay solutions, and nonexistence}
\author{Phuong Le}
\date{Source: arXiv:2609.09724v1 (CC BY 4.0)}
\begin{document}
\maketitle

\noindent\emph{Source and licence.} Radial solutions of quasilinear Hardy--Emden--Fowler equations: exact decay rates, Delaunay solutions, and nonexistence. Version arXiv:2609.09724v1 (CC BY 4.0), distributed by arXiv under the Creative Commons Attribution 4.0 International licence, http://creativecommons.org/licenses/by/4.0/, as recorded in the arXiv licence metadata for this exact version and shown on the abstract page https://arxiv.org/abs/2609.09724v1. Full licence notice: \texttt{licenses/arxiv-2609.09724-CC-BY-4.0.txt}.\par\smallskip

\noindent\textit{Editorial scope.} Counted unit: the article's theorem thm:apriori (source0.tex lines 438--452). The article's complete original proof of it is delivered with it (source0.tex lines 894--970, one \texttt{proof} environment of 77 lines, which the article places away from the statement). No mathematics is written, repaired or re-derived: the statement and the proof are the article's own lines, in the article's own notation, and no retained line is substituted or re-flowed.  Retained uncounted as the source-local context the proof needs: lines 84--87; lines 90--94; lines 348--351; lines 353--357; lines 707--719; lines 857--861.  Nothing was omitted from any retained span, so the package carries no omission record.  No retained line contains a citation, so the wrapper carries no bibliography and no bibliography entry.  No retained line contains a figure environment, an image-inclusion command or drawing code, so no image is delivered and the package declares no asset dependency.  Editorial formatting only, in addition: an AMS \texttt{amsart} preamble that restates the article's own declarations and macros used by the retained lines, namely the environments \texttt{theorem}, \texttt{lemma} on separate counters, together with the standard packages those lines need.  The retained environments print in this document as Theorem 1, Lemma 1 in order of appearance, whereas the article prints the same items with the numbering of its own full document.  The complete native source of the article as distributed in the arXiv e-print for this exact version is bundled unchanged as \texttt{sources/209848/source0.tex}.
\begin{equation}\label{eq:standing}
N>p>1,\qquad s>-p,\qquad 0<\mu<\mub:=\Big(\frac{N-p}{p}\Big)^{p},
\qquad q>p-1,
\end{equation}

\begin{equation}\label{eq:ode}
-\big(r^{N-1}|v'|^{p-2}v'\big)'
=\mu\,r^{N-1-p}v^{p-1}+r^{N-1+s}v^{q},
\qquad r>R_{0}>0 .
\end{equation}

\begin{equation}\label{eq:xY}
x(t):=-\frac{r\,v'(r)}{v(r)},\qquad
Y(t):=r^{p+s}v(r)^{q-p+1},\qquad t=\ln r ,
\end{equation}

\begin{equation}\label{eq:system}
\dot x=\frac{x^{2-p}}{p-1}\big[\mu+Y-h(x)\big],
\qquad
\dot Y=(q-p+1)\,Y\,\big[\gs-x\big] .
\end{equation}

\begin{lemma}\label{lem:system}
Let $v>0$ solve \eqref{eq:ode} on an interval $I\subset(0,\infty)$ with
$v'<0$ on $I$. Then the functions $x,Y$ defined in \eqref{eq:xY} are positive,
of class $C^{1}$ in $t=\ln r$, and satisfy \eqref{eq:system}. Conversely, every
solution $(x,Y)$ of \eqref{eq:system} with $x>0$, $Y>0$ on an interval
$J\subset\R$ arises in this way from a solution $v>0$ of \eqref{eq:ode}, with
$v'<0$, on $\{\mathrm e^{t}:t\in J\}$; $v$ is recovered from
\begin{equation}\label{eq:recover}
v(r)=\big(Y(\ln r)\big)^{\frac{1}{q-p+1}}\,r^{-\gs},
\qquad
\frac{r v'(r)}{v(r)}=-x(\ln r) .
\end{equation}
\end{lemma}

\begin{lemma}\label{lem:auto}
Assume \eqref{eq:standing} and let $v>0$ solve \eqref{eq:ode} on
$(R_{0},+\infty)$. Then $v'(r)<0$ for all large $r$ and
$\lim_{r\to+\infty}v(r)=0$.
\end{lemma}

\begin{theorem}\label{thm:apriori}
Assume \eqref{eq:standing} and let $v>0$ solve \eqref{eq:ode} on
$(R_{0},+\infty)$. Then
\begin{equation}\label{eq:auto}
v'(r)<0\ \text{ for all large }r,
\qquad
\lim_{r\to+\infty}v(r)=0 ,
\end{equation}
and, $(x,Y)$ being as in \eqref{eq:xY}, there is $t_{0}\in\R$ such that
\begin{equation}\label{eq:apriori}
0<x(t)<\nu_{+}\ \text{ for every }t\ge t_{0},
\qquad \liminf_{t\to\infty}x(t)>0,
\qquad \sup_{t\ge t_{0}}Y(t)<+\infty .
\end{equation}
\end{theorem}

\begin{proof}[Proof of Theorem~\ref{thm:apriori}]
Assertion \eqref{eq:auto} is Lemma~\ref{lem:auto}. By that lemma there is
$R_{1}>R_{0}$ with $v'<0$ on $(R_{1},+\infty)$, so Lemma~\ref{lem:system}
applies: $(x,Y)$ solves \eqref{eq:system} on $[t_{0},+\infty)$ with $x>0$,
$Y>0$, where $t_{0}=\ln R_{1}$. It remains to prove \eqref{eq:apriori}.

\smallskip
\noindent\emph{Step 1: $x$ stays away from $0$.} Fix $\delta\in(0,\nu_{-})$;
then $h$ is increasing on $(0,\delta]$ and $h(\delta)<h(\nu_{-})=\mu$. Since
$x^{2-p}>0$ for $x>0$, on $\{0<x\le\delta\}$ we have
\[
\dot x=\frac{x^{2-p}}{p-1}\big[\mu+Y-h(x)\big]
\ \ge\ \frac{x^{2-p}}{p-1}\big(\mu-h(\delta)\big)\ >\ 0 .
\]
Consequently, for every $\eta\in(0,\delta]$ the set $\{x>\eta\}$ is forward
invariant: an orbit reaching $x=\eta$ must cross it upwards. Taking
$\eta:=\min\{x(t_{0}),\delta\}>0$ we obtain $x(t)\ge\eta$ for all $t\ge t_{0}$,
whence $\liminf_{t\to\infty}x(t)\ge\eta>0$.

\smallskip
\noindent\emph{Step 2: $x(t)<\nu_{+}$ for all $t\ge t_{0}$.} We first note that
$x$ cannot blow up in finite time. Indeed $v>0$ and $v\in C^{1}$ on
$(R_{0},+\infty)$, so $x=\frac{r(-v')}{v}$ is continuous on
$(R_{0},+\infty)$ and therefore bounded on every compact subinterval; in
particular $x$ is finite at every finite $t$.

Suppose now, by contradiction, that $x(t_{1})\ge\nu_{+}$ for some
$t_{1}\ge t_{0}$. Since $h$ is strictly decreasing on
$\big(\frac{N-p}{p},\infty\big)$ and $\nu_{+}>\frac{N-p}{p}$, we have
$h(x)\le\mu$ for $x\ge\nu_{+}$, hence
\[
\dot x=\frac{x^{2-p}}{p-1}\big[\mu+Y-h(x)\big]
\ \ge\ \frac{x^{2-p}}{p-1}\,Y\ >\ 0 .
\]
Therefore $x$ is strictly increasing on $[t_{1},\infty)$ and stays above
$\nu_{+}$. Fix $t_{2}>t_{1}$ and put $x_{2}:=x(t_{2})>\nu_{+}$. For
$t\ge t_{2}$,
\[
\dot x\ \ge\ \frac{x^{2-p}}{p-1}\big[\mu-h(x)\big]
=\frac{x^{2-p}}{p-1}\Big[\mu-(N-p)x^{p-1}+(p-1)x^{p}\Big]=:\Theta(x) ,
\]
where $\Theta>0$ on $(\nu_{+},\infty)$ and $\Theta(x)=x^{2}\big(1+o(1)\big)$ as
$x\to+\infty$; in particular
$\int_{x_{2}}^{\infty}\frac{dx}{\Theta(x)}<+\infty$. Comparison with the
scalar equation $\dot z=\Theta(z)$, $z(t_{2})=x_{2}$, then shows that $x$
blows up at a finite time
$t^{*}\le t_{2}+\int_{x_{2}}^{\infty}\frac{dx}{\Theta(x)}$, contradicting the
first paragraph. Hence $x(t)<\nu_{+}$ for all $t\ge t_{0}$.

\smallskip
\noindent\emph{Step 3: $Y$ is bounded.} By Steps 1 and 2 the orbit satisfies
$x(t)\in[\eta,\nu_{+}]$ for all $t\ge t_{0}$, a compact subinterval of
$(0,+\infty)$; hence $\frac{x(t)^{2-p}}{p-1}\ge\vartheta$ for some
$\vartheta>0$. As $h\le\mub$,
on the set $\mathcal Y:=\{t\ge t_{0}:\ Y(t)\ge\mub\}$ we have
$\mu+Y-h(x)\ge\mu>0$, hence
\[
\dot x\ \ge\ \vartheta\,\mu\ =:c_{1}>0\qquad\text{on }\mathcal Y .
\]
Let $J$ be any connected component of $\mathcal Y$ and let $a\ge t_{0}$ be its
left endpoint. Since $x$ increases on $J$ and takes values in
$[\eta,\nu_{+}]$, for every $\tau\in J$ we get
$c_{1}(\tau-a)\le x(\tau)-x(a)\le\nu_{+}-\eta$; hence $J$ is bounded and
\[
|J|\ \le\ T_{1}:=\frac{\nu_{+}-\eta}{c_{1}} .
\]
On the other hand $\dot Y=cY(\gs-x)\le c\,\gs\,Y$ everywhere, so
$Y(t)\le Y(a)e^{c\gs T_{1}}$ for $t\in J$. If $a>t_{0}$ then $Y(a)=\mub$ by
continuity, whence $\sup_{J}Y\le\mub\,e^{c\gs T_{1}}$; if $a=t_{0}$ then
$\sup_{J}Y\le Y(t_{0})e^{c\gs T_{1}}$. Since
$Y<\mub$ outside $\mathcal Y$, we conclude
\[
\sup_{t\ge t_{0}}Y(t)\ \le\
\max\big\{\mub,\ Y(t_{0})\big\}\,e^{c\gs T_{1}}<+\infty .
\]

\end{proof}

\end{document}
